\documentclass[tikz,border=6pt]{standalone}
\usepackage{ctex}
\usepackage{amsmath}
\usepackage{amssymb}
\usepackage{pgfplots}
\pgfplotsset{compat=1.18}
\usetikzlibrary{calc,angles,quotes,arrows.meta,positioning,decorations.markings}
\begin{document}
\begin{tikzpicture}[scale=3.6,>=Stealth]
  % 取 theta = 50 度作为代表点
  \def\ang{50}
  % 坐标轴
  \draw[->,thick] (-1.45,0) -- (1.5,0) node[right] {$\operatorname{Re}$};
  \draw[->,thick] (0,-1.45) -- (0,1.5) node[above] {$\operatorname{Im}$};
  % 单位圆
  \draw[very thick,blue] (0,0) circle (1);
  % 旋转方向箭头（沿圆，表示 theta 增大点逆时针旋转）
  \draw[->,blue!60,line width=1pt]
     (135:1.18) arc (135:200:1.18);
  \node[blue!60,font=\small] at (172:1.34) {$\theta$ 增大};
  % cos theta 投影（实部，红）放在 x 轴上
  \draw[red,line width=1.3pt] (0,0) -- ({cos(\ang)},0);
  % sin theta 投影（虚部，绿）沿 y 轴
  \draw[green!55!black,line width=1.3pt] (0,0) -- (0,{sin(\ang)});
  % 从点到两轴的虚线投影
  \draw[red,dashed]   ({cos(\ang)},{sin(\ang)}) -- ({cos(\ang)},0);
  \draw[green!55!black,dashed] ({cos(\ang)},{sin(\ang)}) -- (0,{sin(\ang)});
  % 半径（从原点到点）
  \draw[very thick,black] (0,0) -- ({cos(\ang)},{sin(\ang)});
  % 辐角弧
  \draw[thick,orange] (0.32,0) arc (0:\ang:0.32);
  \node[orange] at ({0.46*cos(\ang/2)},{0.46*sin(\ang/2)}) {$\theta$};
  % 单位圆上的点
  \filldraw[black] ({cos(\ang)},{sin(\ang)}) circle (0.022);
  \node[above right,font=\small] at ({cos(\ang)},{sin(\ang)})
     {$e^{i\theta}=\cos\theta+i\sin\theta$};
  % 实部、虚部标注
  \node[below,red,font=\small] at ({0.5*cos(\ang)},-0.02) {$\cos\theta$};
  \node[left,green!55!black,font=\small] at (-0.02,{0.5*sin(\ang)}) {$\sin\theta$};
  % 模长 = 1 标注
  \node[font=\small] at ({0.5*cos(\ang)+0.12},{0.5*sin(\ang)+0.14}) {$|e^{i\theta}|=1$};
  % 半透明说明框，置于左下空白处避开主图
  \node[draw=black,fill=yellow!18,opacity=0.9,text opacity=1,
        rounded corners,align=left,font=\footnotesize]
     at (-0.74,-0.85)
     {点在单位圆上\\随 $\theta$ 逆时针旋转\\实部 $=\cos\theta$\\虚部 $=\sin\theta$};
\end{tikzpicture}
\end{document}
